\documentclass[12pt,a4paper]{article}
\usepackage[utf8]{inputenc}
\usepackage{amsmath}
\usepackage{amssymb}
\usepackage{geometry}
\usepackage{xeCJK}
\geometry{margin=1in}

\setCJKmainfont{STSong}

\title{Household-Firm Relationship in HANK Model}
\date{}

\begin{document}

\maketitle

\section{The Basic Structure}

In standard macroeconomic models (including HANK), the relationship is NOT one household per firm.

Instead, we have:
\begin{itemize}
    \item Many households
    \item Many firms
    \item Households own firms (via equity/stocks)
    \item Firms hire workers from households
\end{itemize}

\section{Who Owns Firms?}

In the model, households own firms. This is typically modeled as:

\begin{equation}
\Pi_t = \text{Dividends paid to all households}
\end{equation}

The firm's profits (after paying wages, costs) are distributed as dividends to all households.

\section{How Dividends Work}

\begin{enumerate}
    \item Firms produce output $Y_t$
    \item Firms pay wages $W_t \cdot N_t$
    \item Firms pay other costs
    \item Remaining profit = Dividends = $\Pi_t$
    \item Dividends are distributed to ALL households
\end{enumerate}

Each household receives dividends based on their ownership share.

\section{Three Channels of Income}

A household receives income from:
\begin{enumerate}
    \item Labor Income: $W_t \cdot N_{it}$
    \item Capital Income: $r_t \cdot B_{it}$
    \item Profit Income: $\Pi_{it}$
\end{enumerate}

\section{Key Points}

\begin{enumerate}
    \item Not one-to-one: Households do not each own one firm
    \item Collective ownership: All households collectively own all firms
    \item Three income sources: Labor + Capital + Dividends
    \item Profits flow back: Firm profits $\rightarrow$ Dividends $\rightarrow$ Household income
\end{enumerate}

This is why in the budget constraint we see $\Pi_i$ - it represents the household's share of total firm profits.

\end{document}
